\begin{scanpage}{0009}{v}
In the ``geometric theory of diffraction'' proposed by Keller, the results obtained in the mathematical theory of diffraction of short waves were exactly the ones which were used and generalized. Here, the concept of diffraction rays advanced to the forefront. This concept was expressed rather as a physical hypothesis and was not suitable for representing the field in all of space: it was not usable where the formation of the diffraction field takes place (at the caustic, at the boundary of light and shadow, etc.). Here it is impossible to talk about rays, and one must use a wave interpretation.

What has been said above makes it clear why a large number of works appeared in which the diffraction of short waves was investigated by other methods. Among those applied to reflecting bodies with abrupt surface discontinuities or with sharp edges (strip, disk, finite cylinder or cone, etc.) one should first of all mention the works of P. Ya. Ufimtsev. These works began to appear in print in 1957, and it is on the basis of them that this book was written.

P. Ya. Ufimtsev studied the scattering characteristics by such bodies by taking into account, besides the currents being excited on the surface of the body according to the laws of geometric optics (the ``uniform part of the current'' according to his terminology), the additional currents arising in the vicinity of the edges or borders which have the character of edge waves and rapidly attenuate with increasing distance from the edge or border (the ``nonuniform part of the current''). One may find the radiation field created by the additional currents by comparing the edge or border with the edge of an infinite wedge or the border of a half-plane. In certain cases, one is obliged to consider the diffraction interaction of the various edges --- that is, the fact that the wave created by one edge and propagated past another edge is diffracted by it (secondary diffraction).

Such an approach to the diffraction of short waves has great physical visualizability and allows one to obtain rather simple approximation expressions for the field scattered by various metal bodies. This approach may be called the \uline{physical theory of diffraction}.
\end{scanpage}
